Con todos los datos recopilados y sabiendo que para los algoritmos de ordenamiento de arreglos los deje con un tope de 3600 segundos para decidir el peor no queda claro tomando en cuenta que tanto PatienceSort en general o como QuickSort cuando los arreglos estan ordenados dan los tiempos tope cuando los arreglos son muy grandes. Lo que si podemos hacer es ver el mejor algoritmo que segun los graficos de tiempo para arreglos grandes es el \texttt{std::sort} ya que domina en las 3 distribuciones, aunque si lo vemos para arreglos pequeños el ganador es sin duda QuickSort ya que su tiempo promedio es menor en arreglos mas pequeños. 

Por otro lado al comparar los algoritmos de multiplicacion de matrices nos quedan 2 cosas segun los graficos: la primera es la superioridad del algoritmo Naive en matrices pequeñas ya que como podemos ver en el grafico de tiempo este se iba haciendo mas grande mientras mas grandes sean las matrices. Mientras que para el algoritmo Strassen, por el costo de sus operaciones, va a llegar un momento donde nos es mas optimo utilizar este algoritmo solo cuando las matrices sean muy grandes.